% Section 29.2: The essence of each concept | Part: 7 | Depends on: all previous parts | Status: done | Original lines: 4010-4032
%!TEX root = ../../main.tex
\subsection{The essence of each concept}

{\small
\renewcommand{\arraystretch}{1.3}
\begin{longtable}{@{}p{3.6cm}p{10.4cm}@{}}
\toprule
\textbf{Concept} & \textbf{Essence}\\
\midrule
\endhead
Contraction $\omega(v)=\omega_iv^i$ & canonical pairing $V^*\times V\to\R$; needs no metric. First case $\iota_v\omega$ of the interior product.\\
Exterior product $\alpha\wedge\beta$ & antisymmetrised tensor product, a determinant on $1$-forms; graded commutative; $\dd x$ is the covector dual to $\partial/\partial x$, with $\dd x(\partial_x)=1$.\\
$\int_\gamma\omega$ & contract $\omega$ with the velocity, integrate in $t$: $\int\omega(\dot\gamma)\,\dd t=\int_\gamma\gamma\pb\omega$. Linearity $\Rightarrow$ reparametrisation invariance and sign change with orientation.\\
$\int_\phi\alpha$ & full contraction with the tangent $k$-frame, then Lebesgue integration.\\
Flux $\int_\Sigma\iota_X\mu$ & partial contraction of a volume form with a vector field; volume swept per unit time.\\
Stokes & $\int_M\dd\omega=\int_{\partial M}\omega\iff\dd$ and $\partial$ are adjoint.\\
Gauss & Stokes applied to $\omega=\iota_X\mu$, using $\dd\iota_X\mu=\Lie_X\mu=\operatorname{div}_\mu X\,\mu$.\\
Physics & work $\int F$ is intrinsic (force is a covector); the metric appears only as inertia (kinetic energy), as the vacuum (constitutive laws) or via an observer; fluxes are $2$-forms, charges $3$-forms.\\
Residues & $f\,\dd z$ is closed iff $f$ is holomorphic; $\oint f\,\dd z=2\pi i\sum\Res$ is Stokes on $\Omega$ minus discs. Real part $=$ circulation, imaginary part $=$ flux; the residue is the period over a small circle (de Rham pairing).\\
Geometric algebra & $uv=u\cdot v+u\wedge v$: one product containing the metric and the wedge; $\dd\leftrightarrow\vd\wedge$; one operator $\vd$ for grad, curl, div; $\vd F=J$.\\
Curvature & failure of covariant derivatives to commute: $[\nabla_X,\nabla_Y]-\nabla_{[X,Y]}=R(X,Y)$; the nonzero analogue of $\dd\dd=0$.\\
\bottomrule
\end{longtable}}
